\documentclass[10pt,a4paper]{amsart}
\usepackage[margin=19mm]{geometry}
\usepackage{amsmath,amssymb,mathtools}
\usepackage[scr=rsfs,cal=euler]{mathalfa}
\usepackage{xfrac}
\usepackage[unicode,hidelinks,bookmarks=false]{hyperref}
\newcommand{\ie}{i.e.\ }
\newcommand{\cf}{cf.\ }
\newcommand{\resp}{respectively\ }
\newcommand{\Subsection}{\subsection}
\providecommand{\nobreakdash}{\nobreak\hbox{-}}
\setlength{\emergencystretch}{2em}
\multlinegap0pt
\usepackage{bm}
\usepackage{bbm}
\usepackage{dsfont}
\usepackage{xspace}
\usepackage{cancel}
\usepackage{mathtools}
\usepackage{amsthm}
\makeatletter
\@ifundefined{theorem}{\newtheorem{theorem}{Theorem}}{}
\@ifundefined{lemma}{\newtheorem{lemma}[theorem]{Lemma}}{}
\makeatother
\title[On finite-dimensional homogeneous Lie algebras of derivation]{On finite-dimensional homogeneous Lie algebras of derivations of polynomial rings}
\author{Ivan Arzhantsev, Sergey Gaifullin,, Viktor Lopatkin}
\date{Source: arXiv:2408.05627v2 (CC BY 4.0)}
\begin{document}
\maketitle

\noindent\emph{Source and licence.} On finite-dimensional homogeneous Lie algebras of derivations of polynomial rings. Version arXiv:2408.05627v2 (CC BY 4.0), distributed under the Creative Commons Attribution 4.0 International licence, as recorded in the arXiv licence metadata for this exact version and shown on the abstract page https://arxiv.org/abs/2408.05627v2. Full rights notice: licenses/arxiv-2408.05627-CC-BY-4.0.txt.\par\smallskip

\noindent\textit{Editorial scope.} Counted unit: the Lemma whose source label is \texttt{l4} in the article (source0.tex lines 400--402, source label \texttt{l4}), delivered together with the article's own complete original proof of it (source0.tex lines 404--456, one \texttt{proof} environment of 53 lines). No mathematics is written, repaired or re-derived: statement and proof are the article's own text line for line. Required context retained uncounted: l3 (lines 381--383). Every definition, display and label of those lines is retained unchanged. Omitted from this package and not redistributed: 1--380, 384--399, 403--403, 457--775. Nothing inside a retained excerpt is omitted or altered: every retained body is delivered exactly as the article's lines are written, so no omission receipt of any kind accompanies this package. Citations: the retained excerpts carry no citation command at all, so no bibliography entry of the article is transcribed and no reference is redistributed. Reference adaptation: none was needed and none was made; every reference inside the delivered unit resolves inside it, so no unresolved reference or citation is produced. Printed slips kept literal: no printed word, symbol, hypothesis or number of the counted statement or of its proof was altered. Layout and class: the article's own class is \texttt{amsart} with packages including geometry, babel, fontenc, amsfonts, mathrsfs, tikz, xy, amssymb, amscd, amsthm, amsmath, graphicx, hyperref, array; the wrapper is the lane's pinned \texttt{amsart} editorial wrapper at 10pt on A4, which restates the theorem-like environments the retained text needs and adds the article's own macro definitions that the retained text needs (none), with extra wrapper packages (bm, bbm, dsfont, xspace, cancel, mathtools). The delivered numbering is the unit's own. Assets: no retained span contains a figure environment, a graphics-inclusion command, a PDF-page inclusion, a table or a TikZ picture; no image file is copied into this package, the package declares no asset dependency and no image file is redistributed with it. Licence: version arXiv:2408.05627v2 of this article is distributed under the Creative Commons Attribution 4.0 International licence, as recorded in the arXiv licence metadata for this exact version and shown on the abstract page https://arxiv.org/abs/2408.05627v2. A full rights notice is delivered at licenses/arxiv-2408.05627-CC-BY-4.0.txt. Characters: every retained excerpt is pure ASCII, so the delivered engine, XeLaTeX with its default Latin Modern text font, reports no missing character.
\begin{lemma} \label{l3}
Let $p,q\in\mathbb{Z}^n_{\ge 0}$ be nonzero vectors and the vectors $\beta,\gamma\in\mathbb{K}^n$ are not proportional. If the operators $\Delta^p_{\beta}$ and $\Delta^q_{\gamma}$ generate a finite-dimensional Lie algebra, then there exists $r\in\mathbb{Z}_{\ge 0}$ such that $\langle\beta,rp+q\rangle=0$. 
\end{lemma}

\begin{lemma} \label{l4}
Under the conditions of Lemma~\ref{l3}, let $r$ and $s$ be the smallest non-negative integers such that $\langle\beta,rp+q\rangle=\langle\gamma,p+sq\rangle=0$. Then either $r=0$ or $s=0$. 
\end{lemma}

\begin{proof}
Suppose $r\ge 2$ and $s\ge 1$. Recall that 
$$
[\Delta^q_{\gamma}, \Delta^p_{\beta}]=\Delta^{p+q}_{\langle\gamma,p\rangle\beta-\langle\beta,q\rangle\gamma}.
$$
By assumption, the coefficients at $\beta$ and $\gamma$ in the operator on the right are nonzero. Commuting the operator $\Delta^p_{\beta}$ with $\Delta^q_{\gamma}$ and $\Delta^p_{\beta}$ successively, we obtain
\begin{equation} \label{f1}
[\Delta^q_{\gamma}, \Delta^{tp+(t-1)q}_{u\beta+v\gamma}]=\Delta^{tp+tq}_{\langle\gamma, tp+(t-1)q\rangle(u\beta+v\gamma)-\langle u\beta+v\gamma,q\rangle\gamma}
\end{equation}
and
\begin{equation} \label{f2}
[\Delta^p_{\beta}, \Delta^{tp+tq}_{u'\beta+v'\gamma}]=\Delta^{(t+1)p+tq}_{\langle\beta, tp+tq\rangle(u'\beta+v'\gamma)-\langle u'\beta+v'\gamma,p\rangle\beta}. 
\end{equation}

It follows from the finite dimensionality condition that at some moment these commutators should vanish. Consequently, the coefficients at $\beta$ and $\gamma$ will vanish. 

If we have $\langle\gamma, tp+(t-1)q\rangle u=0$, where $u\ne 0$ and $\langle\gamma,p+sq\rangle=0$, then $1\le s=\frac{t-1}{t}$, and we come to a contradiction. Thus, the last nonzero coefficient at $\beta$ appears in (\ref{f1}).

Now let $\langle\beta,tp+tq\rangle v'=0$ with $v'\ne 0$. Then $\langle\beta,p+q\rangle=0$, which contradicts the assumption $r\ge 2$. Thus, the last nonzero coefficient at $\gamma$ occurs in formula~(\ref{f2}). 

Let us suppose that the coefficient at $\beta$ has become the first to take only zero values. Then, in all of the following steps, the coefficient at $\gamma$ in formula~(\ref{f1}) is equal to
$$
v\langle\gamma,tp+(t-2)q\rangle.
$$ 

Hence, if this coefficient equals zero, we obtain a contradiction because of~${1\le s=\frac{t-2}{t}}$. 

Let us assume now that the coefficient at $\gamma$ has become the first to take only zero values. Then in all subsequent steps, the coefficient at $\beta$ in formula~(\ref{f2}) is equal to
$u'\langle\beta,(t-1)p+tq\rangle.$ If it is equal to zero, then $2\le r=\frac{t-1}{t}$ gives a contradiction. Therefore, neither of these cases is possible.

It remains to consider the case $r=s=1$. We have $\langle\beta,p+q\rangle=\langle\gamma,p+q\rangle=0$, but the pairings 
$$
\langle\beta,p\rangle, \quad \langle\beta,q\rangle, \quad \langle\gamma,p\rangle, \quad \text{and} \quad \langle\gamma,q\rangle 
$$
are nonzero. We obtain
$$
[\Delta^q_{\gamma},[\Delta^q_{\gamma}, \Delta^p_{\beta}]]=[\Delta^q_{\gamma},\Delta^{p+q}_{\langle\gamma,p\rangle\beta-\langle\beta,q\rangle\gamma}]= 
\Delta^{p+2q}_{\omega},
$$
where
$$
\omega=\langle\gamma,p+q\rangle(\langle\gamma,p\rangle\beta-\langle\beta,q\rangle\gamma)-\langle\langle\gamma,p\rangle\beta-\langle\beta,q\rangle\gamma,q\rangle\gamma=v\gamma
$$
with $v=-\langle\langle\gamma,p\rangle\beta-\langle\beta,q\rangle\gamma,q\rangle=\langle\beta,q\rangle\langle\gamma,q-p\rangle\ne 0.$ 

Thus, together with the elements $(\Delta^q_{\gamma}, \Delta^p_{\beta})$, the Lie algebra $\mathfrak{g}(\mathcal{D})$ also contains the elements $(\Delta^{p+2q}_{\gamma}, \Delta^p_{\beta})$ with the condition
$$
\langle\beta,p+2q+p\rangle=\langle\gamma,p+2q+p\rangle=0
$$ 
and $\langle\beta,p\rangle$, $\langle\beta,p+2q\rangle$, $\langle\gamma,p\rangle$, and $\langle\gamma,p+2q\rangle$ are nonzero. 

Repeating the procedure, we pass from this pair to the pair $(\Delta^{p+2(p+2q)}_{\gamma}, \Delta^p_{\beta})$, and so on. This is a contradiction with the finite dimensionality condition. This completes the proof.
\end{proof}

\end{document}
